
% https://akaedu.github.io/book/ch08.html

\chapter{数组}
\BookAnchor{db.ch08}\BookAnchor{db.ch08.id2733202}%


\BookNav{\hyperlink{db.ch07s04}{上一页}\enspace\textperiodcentered\enspace \hyperlink{db.pt01}{上一级}\enspace\textperiodcentered\enspace \hyperlink{db.index}{起始页}\enspace\textperiodcentered\enspace \hyperlink{db.ch08s01}{下一页}}


% https://akaedu.github.io/book/ch08s01.html

\section{数组的基本概念}
\BookAnchor{db.ch08s01}\BookAnchor{db.ch08s01.id2733213}%
数组（Array）\BookAnchor{db.ch08s01.id2732586}也是一种复合数据类型，它由一系列相同类型的元素（Element）\BookAnchor{db.ch08s01.id2732594}组成。例如定义一个由4个\BookCode{int}型元素组成的数组count：


\BookSource{numbers=none}{code/ch08s01/01.pv}{3}{1}


和结构体成员类似，数组\BookCode{count}的4个元素的存储空间也是相邻的。结构体成员可以是基本数据类型，也可以是复合数据类型，数组中的元素也是如此。根据组合规则，我们可以定义一个由4个结构体元素组成的数组：


\BookSource{numbers=none}{code/ch08s01/02.pv}{5}{3}


也可以定义一个包含数组成员的结构体：


\BookSource{numbers=none}{code/ch08s01/03.pv}{6}{4}


数组类型的长度应该用一个整数常量表达式来指定\BookAnchor{db.ch08s01.id2733250}\footnote[16]{\BookAnchor{db.ch08s01.ftn.id2733250}\BookFootnoteBack{db.ch08s01.id2733250}C99的新特性允许在数组长度表达式中使用变量，称为变长数组（VLA，Variable Length Array）\BookAnchor{db.ch08s01.id2733257}，VLA只能定义为局部变量而不能是全局变量，与VLA有关的语法规则比较复杂，而且很多编译器不支持这种新特性，不建议使用。}。数组中的元素通过下标（或者叫索引，Index）\BookAnchor{db.ch08s01.id2733270}来访问。例如前面定义的由4个\BookCode{int}型元素组成的数组\BookCode{count}图示如下：


\BookFigure{数组count}{\BookImage{images/array.count.png}{0.312}}{\BookAnchor{db.ch08s01.id2733290}}


整个数组占了4个\BookCode{int}型的存储单元，存储单元用小方框表示，里面的数字是存储在这个单元中的数据（假设都是0），而框外面的数字是下标，这四个单元分别用\BookCode{count[\BookSplit{}0]\BookSplit{}}、\BookCode{count[\BookSplit{}1]\BookSplit{}}、\BookCode{count[\BookSplit{}2]\BookSplit{}}、\BookCode{count[\BookSplit{}3]\BookSplit{}}来访问。注意，在定义数组\BookCode{int count[\BookSplit{}4]\BookSplit{};\BookSplit{}}时，方括号（Bracket）\BookAnchor{db.ch08s01.id2733345}中的数字4表示数组的长度，而在访问数组时，方括号中的数字表示访问数组的第几个元素。和我们平常数数不同，数组元素是从“第0个”开始数的，大多数编程语言都是这么规定的，所以计算机术语中有Zeroth\BookAnchor{db.ch08s01.id2733360}这个词。这样规定使得访问数组元素非常方便，比如\BookCode{count}数组中的每个元素占4个字节，则\BookCode{count[\BookSplit{}i]\BookSplit{}}表示从数组开头跳过\BookCode{4*i}个字节之后的那个存储单元。这种数组下标的表达式不仅可以表示存储单元中的值，也可以表示存储单元本身，也就是说可以做左值，因此以下语句都是正确的：


\BookSource{numbers=none}{code/ch08s01/04.pv}{5}{3}


到目前为止我们学习了五种后缀运算符：后缀++、后缀--、结构体取成员.、数组取下标[]、函数调用()。还学习了五种单目运算符（或者叫前缀运算符）：前缀++、前缀--、正号+、负号-、逻辑非!。在C语言中后缀运算符的优先级最高，单目运算符的优先级仅次于后缀运算符，比其它运算符的优先级都高，所以上面举例的\BookCode{+\BookSplit{}+\BookSplit{}count[\BookSplit{}2]\BookSplit{}}应该看作对\BookCode{count[\BookSplit{}2]\BookSplit{}}做前缀++运算。



数组下标也可以是表达式，但表达式的值必须是整型的。例如：


\BookSource{numbers=none}{code/ch08s01/05.pv}{4}{2}


使用数组下标不能超出数组的长度范围，这一点在使用变量做数组下标时尤其要注意。C编译器并不检查\BookCode{count[\BookSplit{}-\BookSplit{}1]\BookSplit{}}或是\BookCode{count[\BookSplit{}100]\BookSplit{}}这样的访问越界错误，编译时能顺利通过，所以属于运行时错误\BookAnchor{db.ch08s01.id2733456}\footnote[17]{\BookAnchor{db.ch08s01.ftn.id2733456}\BookFootnoteBack{db.ch08s01.id2733456}你可能会想为什么编译器对这么明显的错误都视而不见？理由一，这种错误并不总是显而易见的，在\hyperlink{db.ch23s01.id2809880}{第 1 节 “指针的基本概念”}会讲到通过指针而不是数组名来访问数组的情况，指针指向数组中的什么位置只有运行时才知道，编译时无法检查是否越界，而运行时每次访问数组元素都检查越界会严重影响性能，所以干脆不检查了；理由二，\hyperlink{db.bi01.bibli.rationale}{[C99 Rationale]}指出C语言的设计精神是：相信每个C程序员都是高手，不要阻止程序员去干他们需要干的事，高手们使用\BookCode{count[\BookSplit{}-\BookSplit{}1]\BookSplit{}}这种技巧其实并不少见，不应该当作错误。}。但有时候这种错误很隐蔽，发生访问越界时程序可能并不会立即崩溃，而执行到后面某个正确的语句时却有可能突然崩溃（在\hyperlink{db.ch10s04.id2742403}{第 4 节 “段错误”}我们会看到这样的例子）。所以从一开始写代码时就要小心避免出问题，事后依靠调试来解决问题的成本是很高的。



数组也可以像结构体一样初始化，未赋初值的元素也是用0来初始化，例如：


\BookSource{numbers=none}{code/ch08s01/06.pv}{3}{1}


则\BookCode{count[\BookSplit{}0]\BookSplit{}}等于3， \BookCode{count[\BookSplit{}1]\BookSplit{}}等于2，后面两个元素等于0。如果定义数组的同时初始化它，也可以不指定数组的长度，例如：


\BookSource{numbers=none}{code/ch08s01/07.pv}{3}{1}


编译器会根据Initializer有三个元素确定数组的长度为3。利用C99的新特性也可以做Memberwise Initialization：


\BookSource{numbers=none}{code/ch08s01/08.pv}{3}{1}


下面举一个完整的例子：


\BookExampleTitle{定义和访问数组}\BookAnchor{db.ch08s01.id2733545}

\BookSource{numbers=none}{code/ch08s01/09.pv}{9}{10}


这个例子通过循环把数组中的每个元素依次访问一遍，在计算机术语中称为遍历（Traversal）\BookAnchor{db.ch08s01.id2733562}。注意控制表达式\BookCode{i \textless{} 4}，如果写成\BookCode{i \textless{}=\BookSplit{} 4}就错了，因为\BookCode{count[\BookSplit{}4]\BookSplit{}}是访问越界。



数组和结构体虽然有很多相似之处，但也有一个显著的不同：数组不能相互赋值或初始化。例如这样是错的：


\BookSource{numbers=none}{code/ch08s01/10.pv}{4}{2}


相互赋值也是错的：


\BookPlain{code/ch08s01/11.txt}{3}{1}


既然不能相互赋值，也就\emph{不能用数组类型作为函数的参数或返回值}。如果写出这样的函数定义：


\BookSource{numbers=none}{code/ch08s01/12.pv}{6}{4}


然后这样调用：


\BookSource{numbers=none}{code/ch08s01/13.pv}{4}{2}


编译器也不会报错，但这样写并不是传一个数组类型参数的意思。对于数组类型有一条特殊规则：\emph{数组类型做右值使用时，自动转换成指向数组首元素的指针}。所以上面的函数调用其实是传一个指针类型的参数，而不是数组类型的参数。接下来的几章里有的函数需要访问数组，我们就把数组定义为全局变量给函数访问，等以后讲了指针再使用传参的办法。这也解释了为什么数组类型不能相互赋值或初始化，例如上面提到的\BookCode{a =\BookSplit{} b}这个表达式，\BookCode{a}和\BookCode{b}都是数组类型的变量，但是\BookCode{b}做右值使用，自动转换成指针类型，而左边仍然是数组类型，所以编译器报的错是\BookCode{error:\BookSplit{} incompatible \BookSplit{}types in assig\BookSplit{}nment}。


\subsection*{习题}
\BookAnchor{db.ch08s01.id2733680}%


1、编写一个程序，定义两个类型和长度都相同的数组，将其中一个数组的所有元素拷贝给另一个。既然数组不能直接赋值，想想应该怎么实现。

\BookNav{\hyperlink{db.ch08}{上一页}\enspace\textperiodcentered\enspace \hyperlink{db.ch08}{上一级}\enspace\textperiodcentered\enspace \hyperlink{db.index}{起始页}\enspace\textperiodcentered\enspace \hyperlink{db.ch08s02}{下一页}}


% https://akaedu.github.io/book/ch08s02.html

\section{数组应用实例：统计随机数}
\BookAnchor{db.ch08s02}\BookAnchor{db.ch08s02.id2733695}%
本节通过一个实例介绍使用数组的一些基本模式。问题是这样的：首先生成一列0\textasciitilde{}\BookSplit{}9的随机数保存在数组中，然后统计其中每个数字出现的次数并打印，检查这些数字的随机性如何。随机数在某些场合（例如游戏程序）是非常有用的，但是用计算机生成完全随机的数却不是那么容易。计算机执行每一条指令的结果都是确定的，没有一条指令产生的是随机数，调用C标准库得到的随机数其实是伪随机（Pseudorandom）\BookAnchor{db.ch08s02.id2733724}数，是用数学公式算出来的确定的数，只不过这些数看起来很随机，并且从统计意义上也很接近均匀分布（Uniform Distribution）\BookAnchor{db.ch08s02.id2733734}的随机数。



C标准库中生成伪随机数的是\BookCode{rand}函数，使用这个函数需要包含头文件\BookCode{stdlib.\BookSplit{}h}，它没有参数，返回值是一个介于0和\BookCode{RAND\_\BookSplit{}MAX}之间的接近均匀分布的整数。\BookCode{RAND\_\BookSplit{}MAX}是该头文件中定义的一个常量，在不同的平台上有不同的取值，但可以肯定它是一个非常大的整数。通常我们用到的随机数是限定在某个范围之中的，例如0\textasciitilde{}\BookSplit{}9，而不是0\textasciitilde{}\BookSplit{}\BookCode{RAND\_\BookSplit{}MAX}，我们可以用\%运算符将\BookCode{rand}函数的返回值处理一下：


\BookSource{numbers=none}{code/ch08s02/01.pv}{3}{1}


完整的程序如下：


\BookExampleTitle{生成并打印随机数}\BookAnchor{db.ch08s02.id2733797}

\BookSource{numbers=left,xleftmargin=19pt}{code/ch08s02/02.pv}{9}{27}


这里介绍一种新的语法：用\BookCode{\#define}定义一个常量。实际上编译器的工作分为两个阶段，先是预处理（Preprocess）\BookAnchor{db.ch08s02.id2733822}阶段，然后才是编译阶段，用\BookCode{gcc}的\BookCode{-\BookSplit{}E}选项可以看到预处理之后、编译之前的程序，例如：


\BookTerminal{code/ch08s02/03.txt}{9}{25}


可见在这里预处理器做了两件事情，一是把头文件\BookCode{stdio.\BookSplit{}h}和\BookCode{stdlib.\BookSplit{}h}在代码中展开，二是把\BookCode{\#define}定义的标识符\BookCode{N}替换成它的定义20（在代码中做了三处替换，分别位于数组的定义中和两个函数中）。像\BookCode{\#include}和\BookCode{\#define}这种以\#号开头的行称为预处理指示（Preprocessing Directive）\BookAnchor{db.ch08s02.id2733895}，我们将在\hyperlink{db.ch21.id2796862}{第 21 章 \emph{预处理}}学习其它预处理指示。此外，用\BookCode{cpp main.\BookSplit{}c}命令也可以达到同样的效果，只做预处理而不编译，\BookCode{cpp}表示C preprocessor。



那么用\BookCode{\#define}定义的常量和\hyperlink{db.ch07s03.id2731288}{第 3 节 “数据类型标志”}讲的枚举常量有什么区别呢？首先，\BookCode{define}不仅用于定义常量，也可以定义更复杂的语法结构，称为宏（Macro）\BookAnchor{db.ch08s02.id2733947}定义。其次，\BookCode{define}定义是在预处理阶段处理的，而枚举是在编译阶段处理的。试试看把\hyperlink{db.ch07s03.id2731288}{第 3 节 “数据类型标志”}习题2的程序改成下面这样是什么结果。


\BookSource{numbers=none}{code/ch08s02/04.pv}{9}{10}


注意，虽然\BookCode{include}和\BookCode{define}在预处理指示中有特殊含义，但它们并不是C语言的关键字，换句话说，它们也可以用作标识符，例如声明\BookCode{int include;\BookSplit{}}或者\BookCode{void define(\BookSplit{}int)\BookSplit{};\BookSplit{}}。在预处理阶段，如果一行以\#号开头，后面跟\BookCode{include}或\BookCode{define}，预处理器就认为这是一条预处理指示，除此之外出现在其它地方的\BookCode{include}或\BookCode{define}预处理器并不关心，只是当成普通标识符交给编译阶段去处理。



回到随机数这个程序继续讨论，一开始为了便于分析和调试，我们取小一点的数组长度，只生成20个随机数，这个程序的运行结果为：


\BookTerminal{code/ch08s02/05.txt}{3}{1}


看起来很随机了。但随机性如何呢？分布得均匀吗？所谓均匀分布，应该每个数出现的概率是一样的。在上面的20个结果中，6出现了5次，而4和8一次也没出现过。但这说明不了什么问题，毕竟我们的样本太少了，才20个数，如果样本足够多，比如说100000个数，统计一下其中每个数字出现的次数也许能说明问题。但总不能把100000个数都打印出来然后挨个去数吧？我们需要写一个函数统计每个数字出现的次数。完整的程序如下：


\BookExampleTitle{统计随机数的分布}\BookAnchor{db.ch08s02.id2734065}

\BookSource{numbers=left,xleftmargin=19pt}{code/ch08s02/06.pv}{9}{33}


我们只要把\BookCode{\#define N}的值改为100000，就相当于把整个程序中所有用到\BookCode{N}的地方都改为100000了。如果我们不这么写，而是在定义数组时直接写成\BookCode{int a[\BookSplit{}20]\BookSplit{};\BookSplit{}}，在每个循环中也直接使用20这个值，这称为硬编码（Hard coding）\BookAnchor{db.ch08s02.id2734105}。如果原来的代码是硬编码的，那么一旦需要把20改成100000就非常麻烦，你需要找遍整个代码，判断哪些20表示这个数组的长度就改为100000，哪些20表示别的数量则不做改动，如果代码很长，这是很容易出错的。所以，\emph{写代码时应尽可能避免硬编码}，这其实也是一个“提取公因式”的过程，和\hyperlink{db.ch07s02.id2730921}{第 2 节 “数据抽象”}讲的抽象具有相同的作用，就是避免一个地方的改动波及到大的范围。这个程序的运行结果如下：


\BookTerminal{code/ch08s02/07.txt}{9}{12}


各数字出现的次数都在10000次左右，可见是比较均匀的。


\subsection*{习题}
\BookAnchor{db.ch08s02.id2734145}%


1、用\BookCode{rand}函数生成[10, 20]之间的随机整数，表达式应该怎么写？

\BookNav{\hyperlink{db.ch08s01}{上一页}\enspace\textperiodcentered\enspace \hyperlink{db.ch08}{上一级}\enspace\textperiodcentered\enspace \hyperlink{db.index}{起始页}\enspace\textperiodcentered\enspace \hyperlink{db.ch08s03}{下一页}}


% https://akaedu.github.io/book/ch08s03.html

\section{数组应用实例：直方图}
\BookAnchor{db.ch08s03}\BookAnchor{db.ch08s03.id2734162}%
继续上面的例子。我们统计一列0\textasciitilde{}\BookSplit{}9的随机数，打印每个数字出现的次数，像这样的统计结果称为直方图（Histogram）\BookAnchor{db.ch08s03.id2734172}。有时候我们并不只是想打印，更想把统计结果保存下来以便做后续处理。我们可以把程序改成这样：


\BookSource{numbers=none}{code/ch08s03/01.pv}{8}{6}


这显然太繁琐了。要是这样的随机数有100个呢？显然这里用数组最合适不过了：


\BookSource{numbers=none}{code/ch08s03/02.pv}{9}{9}


有意思的是，这里的循环变量\BookCode{i}有两个作用，一是作为参数传给\BookCode{howmany}函数，统计数字\BookCode{i}出现的次数，二是做\BookCode{histogram}的下标，也就是“把数字\BookCode{i}出现的次数保存在数组\BookCode{histogram}的第\BookCode{i}个位置”。



尽管上面的方法可以准确地得到统计结果，但是效率很低，这100000个随机数需要从头到尾检查十遍，每一遍检查只统计一种数字的出现次数。其实可以把\BookCode{histogram}中的元素当作累加器来用，这些随机数只需要从头到尾检查一遍（Single Pass）\BookAnchor{db.ch08s03.id2734261}就可以得出结果：


\BookSource{numbers=none}{code/ch08s03/03.pv}{9}{9}


首先把\BookCode{histogram}的所有元素初始化为0，注意使用局部变量的值之前一定要初始化，否则值是不确定的。接下来的代码很有意思，在每次循环中，\BookCode{a[\BookSplit{}i]\BookSplit{}}就是出现的随机数，而这个随机数同时也是\BookCode{histogram}的下标，这个随机数每出现一次就把\BookCode{histogram}中相应的元素加1。



把上面的程序运行几遍，你就会发现每次产生的随机数都是一样的，不仅如此，在别的计算机上运行该程序产生的随机数很可能也是这样的。这正说明了这些数是伪随机数，是用一套确定的公式基于某个初值算出来的，只要初值相同，随后的整个数列就都相同。实际应用中不可能使用每次都一样的随机数，例如开发一个麻将游戏，每次运行这个游戏摸到的牌不应该是一样的。因此，C标准库允许我们自己指定一个初值，然后在此基础上生成伪随机数，这个初值称为Seed\BookAnchor{db.ch08s03.id2734329}，可以用\BookCode{srand}函数指定Seed。通常我们通过别的途径得到一个不确定的数作为Seed，例如调用\BookCode{time}函数得到当前系统时间距1970年1月1日00:\BookSplit{}00:\BookSplit{}00\BookAnchor{db.ch08s03.id2734350}\footnote[18]{\BookAnchor{db.ch08s03.ftn.id2734350}\BookFootnoteBack{db.ch08s03.id2734350}各种派生自UNIX的系统都把这个时刻称为Epoch\BookAnchor{db.ch08s03.id2734354}，因为UNIX系统最早发明于1969年。}的秒数，然后传给\BookCode{srand}：


\BookPlain{code/ch08s03/04.txt}{3}{1}


然后再调用\BookCode{rand}，得到的随机数就和刚才完全不同了。调用\BookCode{time}函数需要包含头文件\BookCode{time.\BookSplit{}h}，这里的\BookCode{NULL}表示空指针，到\hyperlink{db.ch23s01.id2809880}{第 1 节 “指针的基本概念”}再详细解释。


\subsection*{习题}
\BookAnchor{db.ch08s03.id2734410}%


1、补完本节直方图程序的\BookCode{main}函数，以可视化的形式打印直方图。例如上一节统计20个随机数的结果是：


\BookTerminal{code/ch08s03/05.txt}{9}{8}


2、定义一个数组，编程打印它的全排列。比如定义：


\BookSource{numbers=none}{code/ch08s03/06.pv}{4}{2}


则运行结果是：


\BookTerminal{code/ch08s03/07.txt}{9}{8}


程序的主要思路是：


\begin{enumerate}

\item 把第1个数换到最前面来（本来就在最前面），准备打印1xx，再对后两个数2和3做全排列。

\item 把第2个数换到最前面来，准备打印2xx，再对后两个数1和3做全排列。

\item 把第3个数换到最前面来，准备打印3xx，再对后两个数1和2做全排列。
\end{enumerate}


可见这是一个递归的过程，把对整个序列做全排列的问题归结为对它的子序列做全排列的问题，注意我没有描述Base Case怎么处理，你需要自己想。你的程序要具有通用性，如果改变了\BookCode{N}和数组\BookCode{a}的定义（比如改成4个数的数组），其它代码不需要修改就可以做4个数的全排列（共24种排列）。



完成了上述要求之后再考虑第二个问题：如果再定义一个常量\BookCode{M}表示从\BookCode{N}个数中取几个数做排列（\BookCode{N =\BookSplit{}=\BookSplit{} M}时表示全排列），原来的程序应该怎么改？



最后再考虑第三个问题：如果要求从\BookCode{N}个数中取\BookCode{M}个数做组合而不是做排列，就不能用原来的递归过程了，想想组合的递归过程应该怎么描述，编程实现它。

\BookNav{\hyperlink{db.ch08s02}{上一页}\enspace\textperiodcentered\enspace \hyperlink{db.ch08}{上一级}\enspace\textperiodcentered\enspace \hyperlink{db.index}{起始页}\enspace\textperiodcentered\enspace \hyperlink{db.ch08s04}{下一页}}


% https://akaedu.github.io/book/ch08s04.html

\section{字符串}
\BookAnchor{db.ch08s04}\BookAnchor{db.ch08s04.id2734544}%
之前我一直对字符串避而不谈，不做详细解释，现在已经具备了必要的基础知识，可以深入讨论一下字符串了。字符串可以看作一个数组，它的每个元素是字符型的，例如字符串\BookCode{\BookQuoteChar{}Hello,\BookSplit{} world.\BookSplit{}\textbackslash{}n\BookQuoteChar{}}图示如下：


\BookFigure{字符串}{\BookImage{images/array.string.png}{0.638}}{\BookAnchor{db.ch08s04.id2734565}}


注意每个字符末尾都有一个字符\BookCode{'\textbackslash{}0'}做结束符，这里的\BookCode{\textbackslash{}0}是ASCII码的八进制表示，也就是ASCII码为0的Null字符，在C语言中这种字符串也称为以零结尾的字符串（Null-terminated String）\BookAnchor{db.ch08s04.id2734595}。数组元素可以通过数组名加下标的方式访问，而字符串字面值也可以像数组名一样使用，可以加下标访问其中的字符：


\BookSource{numbers=none}{code/ch08s04/01.pv}{3}{1}


但是通过下标修改其中的字符却是不允许的：


\BookPlain{code/ch08s04/02.txt}{3}{1}


这行代码会产生编译错误，说字符串字面值是只读的，不允许修改。字符串字面值还有一点和数组名类似，做右值使用时自动转换成指向首元素的指针，在\hyperlink{db.ch03s03.id2713785}{第 3 节 “形参和实参”}我们看到\BookCode{printf}原型的第一个参数是指针类型，而\BookCode{printf(\BookSplit{}\BookQuoteChar{}hello world\BookQuoteChar{})\BookSplit{}}其实就是传一个指针参数给\BookCode{printf}。



前面讲过数组可以像结构体一样初始化，如果是字符数组，也可以用一个字符串字面值来初始化：


\BookSource{numbers=none}{code/ch08s04/03.pv}{3}{1}


相当于：


\BookSource{numbers=none}{code/ch08s04/04.pv}{3}{1}


\BookCode{str}的后四个元素没有指定，自动初始化为0，即Null字符。注意，虽然字符串字面值\BookCode{\BookQuoteChar{}Hello\BookQuoteChar{}}是只读的，但用它初始化的数组\BookCode{str}却是可读可写的。数组\BookCode{str}中保存了一串字符，以\BookCode{'\textbackslash{}0'}结尾，也可以叫字符串。在本书中只要是以Null字符结尾的一串字符都叫字符串，不管是像\BookCode{str}这样的数组，还是像\BookCode{\BookQuoteChar{}Hello\BookQuoteChar{}}这样的字符串字面值。



如果用于初始化的字符串字面值比数组还长，比如：


\BookSource{numbers=none}{code/ch08s04/05.pv}{3}{1}


则数组\BookCode{str}只包含字符串的前10个字符，不包含Null字符，这种情况编译器会给出警告。如果要用一个字符串字面值准确地初始化一个字符数组，最好的办法是不指定数组的长度，让编译器自己计算：


\BookSource{numbers=none}{code/ch08s04/06.pv}{3}{1}


字符串字面值的长度包括Null字符在内一共15个字符，编译器会确定数组\BookCode{str}的长度为15。



有一种情况需要特别注意，如果用于初始化的字符串字面值比数组刚好长出一个Null字符的长度，比如：


\BookSource{numbers=none}{code/ch08s04/07.pv}{3}{1}


则数组\BookCode{str}不包含Null字符，并且编译器不会给出警告，\hyperlink{db.bi01.bibli.rationale}{[C99 Rationale]}说这样规定是为程序员方便，以前的很多编译器都是这样实现的，不管它有理没理，C标准既然这么规定了我们也没办法，只能自己小心了。



补充一点，\BookCode{printf}函数的格式化字符串中可以用\BookCode{\%s}表示字符串的占位符。在学字符数组以前，我们用\BookCode{\%s}没什么意义，因为


\BookSource{numbers=none}{code/ch08s04/08.pv}{3}{1}


还不如写成


\BookSource{numbers=none}{code/ch08s04/09.pv}{3}{1}


但现在字符串可以保存在一个数组里面，用\BookCode{\%s}来打印就很有必要了：


\BookSource{numbers=none}{code/ch08s04/10.pv}{3}{1}


\BookCode{printf}会从数组\BookCode{str}的开头一直打印到Null字符为止，Null字符本身是Non-printable字符，不打印。这其实是一个危险的信号：如果数组\BookCode{str}中没有Null字符，那么\BookCode{printf}函数就会访问数组越界，后果可能会很诡异：有时候打印出乱码，有时候看起来没错误，有时候引起程序崩溃。

\BookNav{\hyperlink{db.ch08s03}{上一页}\enspace\textperiodcentered\enspace \hyperlink{db.ch08}{上一级}\enspace\textperiodcentered\enspace \hyperlink{db.index}{起始页}\enspace\textperiodcentered\enspace \hyperlink{db.ch08s05}{下一页}}


% https://akaedu.github.io/book/ch08s05.html

\section{多维数组}
\BookAnchor{db.ch08s05}\BookAnchor{db.ch08s05.id2734861}%
就像结构体可以嵌套一样，数组也可以嵌套，一个数组的元素可以是另外一个数组，这样就构成了多维数组（Multi-dimensional Array）\BookAnchor{db.ch08s05.id2734876}。例如定义并初始化一个二维数组：


\BookSource{numbers=none}{code/ch08s05/01.pv}{3}{1}


数组\BookCode{a}有3个元素，\BookCode{a[\BookSplit{}0]\BookSplit{}}、\BookCode{a[\BookSplit{}1]\BookSplit{}}、\BookCode{a[\BookSplit{}2]\BookSplit{}}。每个元素也是一个数组，例如\BookCode{a[\BookSplit{}0]\BookSplit{}}是一个数组，它有两个元素\BookCode{a[\BookSplit{}0]\BookSplit{}[\BookSplit{}0]\BookSplit{}}、\BookCode{a[\BookSplit{}0]\BookSplit{}[\BookSplit{}1]\BookSplit{}}，这两个元素的类型是\BookCode{int}，值分别是1、2，同理，数组\BookCode{a[\BookSplit{}1]\BookSplit{}}的两个元素是3、4，数组\BookCode{a[\BookSplit{}2]\BookSplit{}}的两个元素是5、0。如下图所示：


\BookFigure{多维数组}{\BookImage{images/array.multidim.png}{0.470}}{\BookAnchor{db.ch08s05.id2734954}}


从概念模型上看，这个二维数组是三行两列的表格，元素的两个下标分别是行号和列号。从物理模型上看，这六个元素在存储器中仍然是连续存储的，就像一维数组一样，相当于把概念模型的表格一行一行接起来拼成一串，C语言的这种存储方式称为Row-major\BookAnchor{db.ch08s05.id2734980}方式，而有些编程语言（例如FORTRAN）是把概念模型的表格一列一列接起来拼成一串存储的，称为Column-major\BookAnchor{db.ch08s05.id2734990}方式。



多维数组也可以像嵌套结构体一样用嵌套Initializer初始化，例如上面的二维数组也可以这样初始化：


\BookSource{numbers=none}{code/ch08s05/02.pv}{5}{3}


注意，除了第一维的长度可以由编译器自动计算而不需要指定，其余各维都必须明确指定长度。利用C99的新特性也可以做Memberwise Initialization，例如：


\BookSource{numbers=none}{code/ch08s05/03.pv}{3}{1}


结构体和数组嵌套的情况也可以做Memberwise Initialization，例如：


\BookSource{numbers=none}{code/ch08s05/04.pv}{9}{8}


如果是多维字符数组，也可以嵌套使用字符串字面值做Initializer，例如：


\BookExampleTitle{多维字符数组}\BookAnchor{db.ch08s05.id2735038}

\BookSource{numbers=left,xleftmargin=19pt}{code/ch08s05/05.pv}{9}{18}

\BookFigure{多维字符数组}{\BookImage{images/array.multichar.png}{0.306}}{\BookAnchor{db.ch08s05.id2735044}}


这个程序中定义了一个多维字符数组\BookCode{char days[\BookSplit{}8]\BookSplit{}[\BookSplit{}10]\BookSplit{};\BookSplit{}}，为了使1\textasciitilde{}\BookSplit{}7刚好映射到\BookCode{days[\BookSplit{}1]\BookSplit{}\textasciitilde{}\BookSplit{}days[\BookSplit{}7]\BookSplit{}}，我们把\BookCode{days[\BookSplit{}0]\BookSplit{}}空出来不用，所以第一维的长度是8，为了使最长的字符串\BookCode{\BookQuoteChar{}Wednesday\BookQuoteChar{}}能够保存到一行，末尾还能多出一个Null字符的位置，所以第二维的长度是10。



这个程序和\hyperlink{db.ch04s04.id2719574}{例 4.\BookSplit{}1 “switch语句”}的功能其实是一样的，但是代码简洁多了。简洁的代码不仅可读性强，而且维护成本也低，像\hyperlink{db.ch04s04.id2719574}{例 4.\BookSplit{}1 “switch语句”}那样一堆\BookCode{case}、\BookCode{printf}和\BookCode{break}，如果漏写一个\BookCode{break}就要出Bug。这个程序之所以简洁，是因为用数据代替了代码。具体来说，通过下标访问字符串组成的数组可以代替一堆\BookCode{case}分支判断，这样就可以把每个\BookCode{case}里重复的代码（\BookCode{printf}调用）提取出来，从而又一次达到了“提取公因式”的效果。这种方法称为数据驱动的编程（Data-driven Programming）\BookAnchor{db.ch08s05.id2735162}，写代码最重要的是选择正确的数据结构来组织信息，设计控制流程和算法尚在其次，只要数据结构选择得正确，其它代码自然而然就变得容易理解和维护了，就像这里的\BookCode{printf}自然而然就被提取出来了。\hyperlink{db.bi01.bibli.manmonth}{[人月神话]}中说过：“Show me your flowcharts and conceal your tables, and I shall continue to be mystified.\BookSplit{} Show me your tables, and I won't usually need your flowcharts; they'll be obvious.\BookSplit{}”



最后，综合本章的知识，我们来写一个最简单的小游戏－－剪刀石头布：


\BookExampleTitle{剪刀石头布}\BookAnchor{db.ch08s05.id2735200}

\BookSource{numbers=left,xleftmargin=19pt}{code/ch08s05/06.pv}{9}{31}


0、1、2三个整数分别是剪刀石头布在程序中的内部表示，用户也要求输入0、1或2，然后和计算机随机生成的0、1或2比胜负。这个程序的主体是一个死循环，需要按Ctrl-C退出程序。以往我们写的程序都只有打印输出，在这个程序中我们第一次碰到处理用户输入的情况。我们简单介绍一下\BookCode{scanf}函数的用法，到\hyperlink{db.ch25s02.id2832755}{第 2.\BookSplit{}9 节 “格式化I/\BookSplit{}O函数”}再详细解释。\BookCode{scanf(\BookSplit{}\BookQuoteChar{}\%d\BookQuoteChar{},\BookSplit{} \&man)\BookSplit{}}这个调用的功能是等待用户输入一个整数并回车，这个整数会被\BookCode{scanf}函数保存在\BookCode{man}这个整型变量里。如果用户输入合法（输入的确实是数字而不是别的字符），则\BookCode{scanf}函数返回1，表示成功读入一个数据。但即使用户输入的是整数，我们还需要进一步检查是不是在0\textasciitilde{}\BookSplit{}2的范围内，写程序时对用户输入要格外小心，用户有可能输入任何数据，他才不管游戏规则是什么。



和\BookCode{printf}类似，\BookCode{scanf}也可以用\BookCode{\%c}、\BookCode{\%f}、\BookCode{\%s}等转换说明。如果在传给\BookCode{scanf}的第一个参数中用\BookCode{\%d}、\BookCode{\%f}或\BookCode{\%c}表示读入一个整数、浮点数或字符，则第二个参数的形式应该是\&运算符加相应类型的变量名，表示读进来的数保存到这个变量中，\&运算符的作用是得到一个指针类型，到\hyperlink{db.ch23s01.id2809880}{第 1 节 “指针的基本概念”}再详细解释；如果在第一个参数中用\BookCode{\%s}读入一个字符串，则第二个参数应该是数组名，数组名前面不加\&，因为数组类型做右值时自动转换成指针类型，在\hyperlink{db.ch10s02.id2741278}{第 2 节 “断点”}有\BookCode{scanf}读入字符串的例子。



留给读者思考的问题是：\BookCode{(\BookSplit{}man -\BookSplit{} computer +\BookSplit{} 4)\BookSplit{} \% 3 -\BookSplit{} 1}这个神奇的表达式是如何比较出0、1、2这三个数字在“剪刀石头布”意义上的大小的？

\BookNav{\hyperlink{db.ch08s04}{上一页}\enspace\textperiodcentered\enspace \hyperlink{db.ch08}{上一级}\enspace\textperiodcentered\enspace \hyperlink{db.index}{起始页}\enspace\textperiodcentered\enspace \hyperlink{db.ch09}{下一页}}
